\chapter*{TÓM TẮT}
\addcontentsline{toc}{chapter}{TÓM TẮT}

Khóa luận \emph{\ThesisTitle} nghiên cứu dự báo tiến triển MCI dưới dạng bài toán sống còn dọc đa phương thức. Từ khung CNN--T-LSTM của Aghajanian và cộng sự, nghiên cứu phát triển lần lượt các baseline MRI, ghép nối đa phương thức, tương tác từng cặp và kiến trúc hợp nhất gọn hơn; phân rã tensor được đưa vào sau khi đánh giá chi phí của outer-product.

Cohort ADNI đã đóng băng gồm 359 người, với 106 biến cố. Ba MRI được chọn trong cửa sổ đủ điều kiện 18 tháng. MRI thứ ba xác định mốc tham chiếu dự báo; biến cố được xác định sau landmark đủ điều kiện bằng ngưỡng MMSE/CDR. Dữ liệu clinical được ghép theo ngày gần nhất trong cửa sổ $\pm90$ ngày quanh MRI, nên tính sẵn có của toàn bộ đầu vào tại mốc dự báo cần được kiểm tra riêng khi triển khai tiến cứu.

Ba tương tác MRI--Radiomics (MR), MRI--Clinical (MC) và Radiomics--Clinical (RC) ban đầu dùng outer-product và định tuyến residual có cổng động. CP sau đó tham số hóa tensor \emph{trọng số} của phép song tuyến, giữ nguyên MLP phía sau. Lõi tương tác ba cặp giảm từ 229,696 xuống 3,712 tham số; validation C-index trung bình ba seed thay đổi từ 0.8344 xuống 0.8278. Ablation dẫn đến kiến trúc RC-Free Pairwise Fusion: vẫn giữ cả ba modality nhưng chỉ duy trì MR/MC. Bốn route có sigmoid gate phụ thuộc đầu vào và scale toàn cục học được, trước T-LSTM, attention theo thời gian và Cox head.

Trên benchmark chuẩn hóa, RC-Free Pairwise Fusion đạt validation C-index $0.8347\pm0.0237$ (độ lệch chuẩn mẫu qua ba seed). Kết quả ủng hộ một kiến trúc gọn có hiệu năng cạnh tranh, chưa chứng minh ưu thế thống kê. Baseline sống còn cổ điển vẫn mạnh theo quy trình lựa chọn và quy ước concordance riêng. Kiến trúc được chọn bằng validation; việc truy cập test trong lịch sử được ghi nhận như một giới hạn. Đánh giá cuối cùng theo protocol held-out chính thức chưa được thực hiện.

\textbf{Từ khóa:} Alzheimer, MCI, MRI dọc, radiomics, hợp nhất tensor, phân rã CP, gated residual, T-LSTM, phân tích sống còn.
